% 07_ketluan.tex - Phần 6: Kết luận
\section{Kết luận}

\subsection{Tổng kết}

Dự án đã giải quyết trọn vẹn yêu cầu của đề tài thông qua việc tái định nghĩa các vấn đề xã hội bằng ngôn ngữ kỹ thuật:

\begin{enumerate}
    \item \textbf{Định lượng hóa thành công:} Chuyển đổi "nghiện" thành RPE, "điều hướng" thành drift velocity, "rủi ro" thành risk index.
    
    \item \textbf{Giải pháp tự động:} Xây dựng Algorithmic Circuit Breaker dựa trên PID control, hoạt động liên tục mà không cần giám sát thủ công.
    
    \item \textbf{Chứng minh khả thi:} Mô phỏng cho kết quả giảm 64.6\% tỷ lệ kiệt sức và tăng 15.4\% retention, chứng tỏ tính hiệu quả.
    
    \item \textbf{Phân tích toàn diện:} Xem xét đầy đủ các góc độ kỹ thuật, đạo đức, pháp lý và kinh tế.
\end{enumerate}

\textbf{Algorithmic Circuit Breaker không chỉ là một tính năng, mà là một kiến trúc an toàn cần thiết cho mọi hệ thống trí tuệ nhân tạo hiện đại.}

\subsection{Thông điệp của nhóm kỹ sư}

Chúng tôi khẳng định rằng:

\begin{quote}
\textit{"Trách nhiệm của kỹ sư không phải là giáo dục hành vi cho người dùng, mà là đảm bảo hệ thống do mình tạo ra không vận hành vượt quá giới hạn an toàn của con người. Đó chính là cốt lõi của đạo đức nghề nghiệp trong kỷ nguyên số."}
\end{quote}

Tương tự như kỹ sư cơ khí phải đảm bảo máy móc có circuit breaker điện, kỹ sư phần mềm phải đảm bảo thuật toán có circuit breaker tâm lý.

\subsection{Đóng góp chính của đề tài}

\begin{enumerate}
    \item \textbf{Đóng góp lý thuyết:}
    \begin{itemize}
        \item Mô hình hóa tác động tâm lý bằng lý thuyết điều khiển
        \item Định nghĩa drift velocity trong không gian vector nhúng
        \item Chứng minh mối quan hệ giữa bảo vệ người dùng và $CLV$
    \end{itemize}
    
    \item \textbf{Đóng góp thực tiễn:}
    \begin{itemize}
        \item Kiến trúc middleware có thể tích hợp vào hệ thống hiện có
        \item Code mô phỏng mã nguồn mở (sẽ công khai trên GitHub)
        \item Đề xuất metrics cụ thể để đo lường well-being
    \end{itemize}
    
    \item \textbf{Đóng góp xã hội:}
    \begin{itemize}
        \item Nâng cao nhận thức về trách nhiệm của kỹ sư AI
        \item Cung cấp framework cho regulators và policymakers
        \item Bảo vệ sức khỏe tâm thần của hàng tỷ người dùng
    \end{itemize}
\end{enumerate}

\subsection{Hạn chế và hướng phát triển}

\textbf{Hạn chế:}
\begin{itemize}
    \item Chưa thử nghiệm trên người dùng thực
    \item Mô hình hóa đơn giản chưa phản ánh đầy đủ độ phức tạp của tâm lý con người
    \item Chưa giải quyết vấn đề privacy khi thu thập dữ liệu chi tiết
\end{itemize}

\textbf{Hướng phát triển:}
\begin{itemize}
    \item Thử nghiệm lâm sàng với IRB approval
    \item Mở rộng sang các nền tảng khác: gaming, e-commerce
    \item Chuẩn hóa thành IEEE/ISO standard
    \item Xây dựng open-source toolkit cho developer
\end{itemize}

\subsection{Lời kết}

Mạng xã hội đã trở thành một phần không thể thiếu của đời sống hiện đại. Nhưng như mọi công nghệ mạnh mẽ khác, nó cần được kiểm soát để phục vụ con người, không phải ngược lại.

Algorithmic Circuit Breaker là bước đầu tiên trong hành trình biến AI từ công cụ tối đa hóa lợi nhuận thành công cụ bảo vệ sức khỏe con người. Chúng tôi hy vọng nghiên cứu này sẽ truyền cảm hứng cho các kỹ sư, nhà nghiên cứu và nhà hoạch định chính sách cùng nhau xây dựng một không gian số an toàn và bền vững hơn.

\textbf{Tương lai của công nghệ không chỉ phụ thuộc vào khả năng kỹ thuật, mà còn phụ thuộc vào trách nhiệm đạo đức của những người tạo ra nó.}

\newpage
